\begin{table}[H]
\centering
\captionsetup{justification=centering}
\caption{Hiperparâmetros estimados para o total de ocupados - modelos univariado e multivariado}
\label{tab:hiperocup}
\scalebox{0.92}{
\renewcommand{\arraystretch}{0.85}
\begin{tabular}{llccccc}
\toprule
& & \multicolumn{5}{c}{\textbf{Modelo Univariado}} \\
\cmidrule(lr){3-7}
\textbf{Estrato geográfico} & \textbf{Processo} & \(\hat{\sigma}_L^2\) & \(\hat{\sigma}_R^2\) & \(\hat{\sigma}_S^2\) & \(\hat{\sigma}_I^2\) & \(\hat{\sigma}_{\tilde{e}}^2\) \\
\midrule
01 - Belo Horizonte & AR(1) & 429,6609 & 0,8061 & 0,0000 & 0,0000 & 0,5640 \\
02 - Entorno e Colar Metropolitano de BH & AR(1) & 826,3711 & 6,7883 & 0,0000 & 0,0000 & 0,5534 \\
03 - Sul de Minas & AR(1) & 0,3494 & 4,4507 & 0,0000 & 0,0000 & 0,4652 \\
04 - Triângulo Mineiro & AR(1) & 7,8799 & 7,8544 & 0,0000 & 0,0000 & 0,3975 \\
05 - Zona da Mata & AR(1) & 5,1703 & 10,4187 & 0,0000 & 0,0000 & 0,5184 \\
06 - Norte de Minas & AR(1) & 607,6945 & 2,3814 & 0,0000 & 0,0023 & 0,5275 \\
07 - Vale do Rio Doce & AR(1) & 0,0000 & 0,0005 & 0,0022 & 0,0001 & 0,3867 \\
08 - Central & AR(1) & 531,9691 & 1,2636 & 0,0000 & 0,0000 & 0,5782 \\
\midrule
& & \multicolumn{5}{c}{\textbf{Modelo Multivariado}} \\
\cmidrule(lr){3-7}
\textbf{Estrato geográfico} & \textbf{Processo} & \(\hat{\sigma}_L^2\) & \(\hat{\sigma}_R^2\) & \(\hat{\sigma}_S^2\) & \(\hat{\sigma}_I^2\) & \(\hat{\sigma}_{\tilde{e}}^2\) \\
\midrule
01 - Belo Horizonte & AR(1) & 429,6600 & 0,8061 & 0,0000 & 0,0000 & 0,5640 \\
02 - Entorno e Colar Metropolitano de BH & AR(1) & 826,3719 & 7,0885 & 0,0000 & 0,0000 & 0,5534 \\
03 - Sul de Minas & AR(1) & 0,3494 & 5,0512 & 0,0000 & 0,0000 & 0,4652 \\
04 - Triângulo Mineiro & AR(1) & 7,8798 & 8,7550 & 0,0000 & 0,0000 & 0,3975 \\
05 - Zona da Mata & AR(1) & 5,1704 & 11,6196 & 0,0000 & 0,0000 & 0,5184 \\
06 - Norte de Minas & AR(1) & 607,6952 & 3,8825 & 0,0000 & 0,0023 & 0,5275 \\
07 - Vale do Rio Doce & AR(1) & 0,0000 & 1,8018 & 0,0022 & 0,0001 & 0,3867 \\
08 - Central & AR(1) & 531,9696 & 3,3652 & 0,0000 & 0,0000 & 0,5782 \\
\bottomrule
\end{tabular}}
\fonte{Elaboração própria, com base nos dados da PNAD Contínua. Nota: o processo do erro amostral é identificado individualmente por estrato (Seção \ref{sec:metodologia}). Impõe-se \(V(\tilde{e}_t)=1\), de modo que a estimativa direta preserva a variância do desenho; a variância da inovação \(\hat{\sigma}_{\tilde{e}}^2\) não é estimada, mas \textbf{derivada} dos coeficientes do processo pela equação de Lyapunov, e por isso coincide nos dois modelos. O componente irregular é fracamente identificado: o perfil de verossimilhança é quase plano em \(\sigma_I^2\), e suas estimativas devem ser lidas com cautela.}
\end{table}
